% appB-offline.tex -- full deferred proofs of Section 3.
\section{Full proof of the offline characterization}\label{app:offline}

Throughout, $A\in\{0,1\}^{n\times T}$ is the incidence matrix ($A_{it}=\1[i\in\supp(v_t)]$), $r_i(S)=\sum_{t\in S}A_{it}$,
$\mu_s=ds/n$, $\ell_s=\log(eT/s)$. We assume $\tau_-\,n\le T\le\tau_+\,n$ for absolute constants $0<\tau_-\le\tau_+$, so
$\ell_s=\Th(\log(n/s))$ and $\ell_s\ge\log e=1$.

\subsection{The lower bound (\Cref{lem:offline-lower})}\label{app:offline-lower-full}

\paragraph{Poisson tails and the inversion.}
For $X\sim\Pois(\mu)$, the Cram\'er rate at $\mu+t$ is $I(t)=\mu\,h(t/\mu)$ with $h(u)=(1+u)\log(1+u)-u$. Since
$h(u)=\Th(u^2)$ for $u\le 1$ and $h(u)=\Th(u\log(1+u))$ for $u\ge 1$, we get \eqref{eq:poisson-tail}. Solving rate $=\Lscale$:
in the regime $t\lesssim\mu$, $t^2/\mu\asymp\Lscale$ gives $t\asymp\sqrt{\mu\Lscale}$; in the regime $t\gtrsim\mu$, setting
$t=c\,\Lscale/\log(e+\Lscale/\mu)$ and $z=\Lscale/\mu\ge 1$ gives $t\log(1+t/\mu)\asymp \mu\cdot\frac{z}{\log(e+z)}\cdot\log\!\big(1+\tfrac{z}{\log(e+z)}\big)\asymp\mu z=\Lscale$, because $\log(1+z/\log(e+z))\asymp\log(e+z)$ for $z\ge 1$. Hence the deviation at which a single $\Pois(\mu)$ load reaches rate $\Lscale$ is $t^*(\mu,\Lscale)=\Th(\sqrt{\mu\Lscale}+\Lscale/\log(e+\Lscale/\mu))=\Th(\Psifn(\mu,\Lscale))$.

\paragraph{Step 1: means.}
Fix a coloring with means $\mu_c=m_c d/n$, $\sum_c\mu_c=k\lam$. The loads are independent with $\E Y_c=\mu_c$ and
$\Var Y_c=\mu_c(1-d/n)=\Th(\mu_c)$. We do not assume balance; instead we lower-bound the range by a dichotomy on the means.
Throughout, $t^*=t^*(\lam,\log k)=\Th(\Psifn(\lam,\log k))$ is the deviation of the Poisson inversion above, and
$t^*\ge 4\sqrt{\lam\log k}$ (since $\Psifn\ge\sqrt{\lam\log k}$ and $\log k$ is large).

\paragraph{Step 2: Poissonization.}
Let $B$ be the $n\times T$ matrix with i.i.d.\ $\Bin(1,p)$ entries, $p=d/n$. Conditioning each column of $B$ on having
exactly $d$ ones yields $\mu_{d,n}$, and by Stirling $\Pr(\Bin(n,p)=d)\ge c_0/\sqrt{d+1}$, so $\Pr(\mathcal C)\ge(c_0/\sqrt{d+1})^T$
where $\mathcal C$ is the event that all columns have weight $d$. Hence for any event $E$,
$\Pr_{\mathrm{exact}}(E)\le\Pr_B(E)/\Pr_B(\mathcal C)\le(C\sqrt{d+1})^T\Pr_B(E)$. Under $B$ the $k$ loads at a fixed
coordinate, $Y_c=\sum_{t\in\chi^{-1}(c)}B_{it}\sim\Bin(m_c,p)$, are independent, with $\E Y_c=\mu_c=m_c p$ and
$\Var Y_c=\Th(\mu_c)$ and $\sum_c\mu_c=k\lam$.

\paragraph{Step 3: per-coordinate spread.}
Fix a coordinate. The $k$ loads are independent with means $\mu_c$ summing to $k\lam$, and the range $\rangeop_i=\max_c Y_c-\min_c Y_c$
is the max-minus-min load of a weighted balls-into-bins allocation of (in expectation) $k\lam$ balls into $k$ bins. Write
$t^*=\Th(\Psifn(\lam,\log k))$ for the Poisson deviation of the inversion above. We use the following lower tail of the
occupancy spread.

\begin{lemma}[Occupancy spread, lower tail]\label{lem:occ-spread}
Let $Y_1,\dots,Y_k$ be independent with $Y_c\sim\Bin(m_c,p)$, $\mu_c=m_c p$, $\sum_c\mu_c=k\lam$, and $\lam=\Th(d/k)$. Then
$\Pr\big(\max_c Y_c-\min_c Y_c\le\tfrac14 t^*\big)\le e^{-q}$ with $q\ge k^{1/5}$.
\end{lemma}

The proof rests on one standard occupancy input, the comparable-mean extreme-value law: for $m$ independent loads of
\emph{comparable} mean $\Th(M)$ (means within a factor $2$), the maximum-minus-minimum is $\ge\Th(\Psifn(M,\log m))$ off an
event of probability $e^{-\Th(\sqrt m)}$ \cite{ABKU1999,BCSV2006}. Quantitatively: a single load exceeds the common level $M$
by $\Om(\Psifn(M,\log m))$ with probability $\ge m^{-1/2}$ (rate $\tfrac12\log m$), so the maximum does unless all $m$ loads
fail, probability $\le(1-m^{-1/2})^m\le e^{-\sqrt m}$; and the minimum is $\le M$ \whp\ (the median load is $M+\Oh(1)$, so
$\Pr(\min>M)\le 2^{-m}$). The difference is then $\ge\Om(\Psifn(M,\log m))$. (The matching
\emph{typical} value
$\rangeop_i=\Th(\Psifn(\lam,\log k))$ \whp\ is the weighted maximum-load law from the same references; \Cref{lem:occ-spread} is
the lower-tail refinement we need.) We extract a large comparable-mean sub-family and apply this law.

\emph{Heavy bin.} If some $\mu_{c^\dagger}\ge k^{1/4}\lam$, the average constraint $\sum_c\mu_c=k\lam$ forces $\ge\Om(k^{1/4})$
bins of mean $\le\lam$ (each heavy bin consumes mean $\ge k^{1/4}\lam$). The heavy bin has $Y_{c^\dagger}\ge\mu_{c^\dagger}/2\ge\tfrac12 k^{1/4}\lam$
by a Chernoff bound (failure $e^{-\Om(k^{1/4}\lam)}$). Each sub-$\lam$ bin (mean $\le\lam$) is $\le 2\lam$ with probability
$\ge\tfrac12$ (Markov), so the minimum over the $\Om(k^{1/4})$ independent sub-$\lam$ bins exceeds $2\lam$ with probability
$\le 2^{-\Om(k^{1/4})}$. Since $k^{1/4}\lam\ge 8(\lam+t^*)$ (as $t^*\le\Psifn(\lam,\log k)=(\log n)^{o(1)}$ and $k^{1/4}\to\infty$),
$\rangeop_i\ge\tfrac12 k^{1/4}\lam-2\lam\ge\tfrac14 k^{1/4}\lam\ge t^*$ \whp, with $q=\Om(k^{1/4})$.

\emph{No heavy bin.} Otherwise every $\mu_c<k^{1/4}\lam$. Discard bins with $\mu_c<\lam/2$ (they only lower the minimum); the
rest $H=\{c:\lam/2\le\mu_c<k^{1/4}\lam\}$ carry $\sum_{c\in H}\mu_c>k\lam/2$, so $|H|\ge(k\lam/2)/(k^{1/4}\lam)=k^{3/4}/2$.
Splitting $[\lam/2,k^{1/4}\lam)$ into $\Oh(\log k)$ dyadic bands, the densest band $B$ has $k'':=|B|\ge|H|/\Oh(\log k)\ge k^{3/4-o(1)}$
bins of comparable mean $M\in[\lam/2,k^{1/4}\lam)$. Applying the comparable-mean law to $B$ (with $m=k''$, $\log k''=\Th(\log k)$):
$\max_{c\in B}Y_c-\min_{c\in B}Y_c\ge c_1\Psifn(M,\log k'')\ge c_1\Psifn(\lam/2,\Th(\log k))\ge\tfrac14 t^*$ \whp, since
$\Psifn$ is nondecreasing in its first argument, $M\ge\lam/2$, and $\Psifn(\lam/2,\Th(\log k))\ge c_2\Psifn(\lam,\log k)=c_2 t^*$
for absolute $c_1,c_2$ (halving $\lam$ and the logarithm changes $\Psifn$ by a constant factor); the constants are absorbed
into the threshold $\tfrac14 t^*$. The failure probability is $e^{-\Th(\sqrt{k''})}=e^{-k^{3/8-o(1)}}$, so $q\ge k^{1/5}$.
This proves \Cref{lem:occ-spread}.

Step~2 placed us in the Bernoulli model $B$, whose rows are independent; the events $\{\rangeop_i\le D\}_{i\in[n]}$ depend on
disjoint rows of $B$, so they are independent and $\Pr_B(\max_i\rangeop_i\le\tfrac14 t^*)\le e^{-qn}$ by \Cref{lem:occ-spread}.
Set $D=\tfrac14 t^*=\Th(\Psifn)$.

\paragraph{Step 4: union over colorings.}
Step 3 gives a per-coordinate rate $q\ge k^{1/5}$. Combining Steps 2--3 and union-bounding over the at most $k^T$ colorings,
\[
  \Pr_{\mathrm{exact}}\big(\exists\chi:\ \max_i\rangeop_i(T)\le D\big)\ \le\ (C\sqrt{d+1})^T\,k^T\,e^{-q n}
  \ =\ \exp\!\big(\Oh(n\log(dk))-\Om(n\,k^{1/5})\big).
\]
By the hypothesis $k\ge C_0\log^5(dk)$, $k^{1/5}\ge C_0^{1/5}\log(dk)$ dominates the $\Oh(\log(dk))$ per-coordinate cost, so
the probability is $\exp(-\Om(n\,k^{1/5}))=o(1)$. Therefore $\offdisc_k>D=\tfrac14 t^*\ge c\,\Psifn(\lam,\log k)$ \whp. The three
regimes \eqref{eq:Psi-regimes} are the three behaviours of $t^*(\lam,\log k)$: Gaussian $\sqrt{\lam\log k}$ for
$\lam\gtrsim\log k$, the Bennett term $\log k/\log((\log k)/\lam)$ for $1\ll\lam\ll\log k$, and $\log k/\loglog k$ for
$\lam=\Th(1)$. \qed

\subsection{The structural lemma (\Cref{lem:structural})}\label{app:structural-full}

The structural lemma asserts that, simultaneously for every subset $S$ of columns, the row loads $r_i(S)$ are not too
concentrated on any small set of rows. The proof has the shape one expects for a statement of this kind. We first show that
a single row is unlikely to be heavily loaded, controlling the expected number of overloaded rows by a moment computation;
we then upgrade this expectation bound to a high-probability statement using a tail inequality that survives the dependence
between rows; finally we pay for the union over all $\binom{T}{s}$ choices of $S$ and over all $s$. The one structural
ingredient that makes the upgrade possible is negative association, which lets us treat the row loads as if they were
independent for the purpose of upper-tail bounds.

\paragraph{Negative association.}
We begin by recording why the row loads are negatively associated, since this is what licenses the Chernoff-type
factorization used below. The indicator vector of a uniform $d$-subset of $[n]$ is strongly Rayleigh, and strong Rayleigh
measures are negatively associated across coordinates \cite{BorceaBrandenLiggett2009}. The $s$ columns indexed by $S$ are
drawn independently, so concatenating their indicator vectors gives a negatively associated vector on the product index set
$\{(i,t):i\in[n],t\in S\}$; independence of the blocks preserves NA. Now each row load $r_i(S)=\sum_{t\in S}A_{it}$ is a
nondecreasing function of the block $\{(i,t):t\in S\}$, and distinct rows $i$ depend on disjoint blocks. Negative
association is closed under taking nondecreasing functions of disjoint blocks of variables, so the family
$\{r_i(S)\}_{i\in[n]}$ is itself NA (\Cref{app:prelims}). This is exactly the property we need: for nondecreasing functions
of the $r_i$, the upper tail of a sum factorizes as if the rows were independent.

\paragraph{The two tail estimates.}
Fix $S$ and write $s=|S|$. Each row load $r_i(S)$ is a sum of $s$ independent $\Bin(1,d/n)$ indicators, so
$r_i(S)\sim\Bin(s,d/n)$ with mean $\mu_s=ds/n$. Fix a large constant $B$; it will be tuned at the very end so that the union
bound closes. The level above which a row counts as overloaded changes character depending on whether the mean $\mu_s$ is
above or below the logarithmic budget $\ell_s$, so we treat the two regimes separately. In each case we exhibit a
nonnegative variable $Y_i$, a nondecreasing function of $r_i$, whose expectation is at most $e^{-(B-1)\ell_s}$; this is the
per-row input the high-probability step consumes.

\emph{Dense case $\mu_s\ge\ell_s$.} Here the natural overload scale is the multiplicative one, so we use a soft exponential
test variable in place of a hard threshold. Let $Y_i=\exp(-K^2\mu_s\ell_s/(16\,r_i))\1[r_i>0]$, which is increasing in
$r_i$ and bounded by $1$. Split on whether $r_i$ exceeds $b\mu_s$ for a constant $b$ to be chosen. On the event
$\{0<r_i\le b\mu_s\}$ the exponent $-K^2\mu_s\ell_s/(16 r_i)$ is at most $-K^2\mu_s\ell_s/(16 b\mu_s)=-K^2\ell_s/(16b)$, so
$Y_i\le e^{-K^2\ell_s/(16 b)}$ there. On the complementary event $\{r_i>b\mu_s\}$ we only use $Y_i\le 1$, but this event is
itself rare: a multiplicative Chernoff bound for $\Bin(s,d/n)$ gives $\Pr(r_i>b\mu_s)\le e^{-c_b\mu_s}\le e^{-c_b\ell_s}$,
the last step using $\mu_s\ge\ell_s$. Combining the two events, $\E Y_i\le e^{-K^2\ell_s/(16b)}+e^{-c_b\ell_s}$. We now fix
constants in order: choose $b$ large enough that $c_b\ge B$, and then choose $K$ large enough that $K^2/(16b)\ge B$. Both
summands are then at most $e^{-B\ell_s}$, so $\E Y_i\le 2e^{-B\ell_s}\le e^{-(B-1)\ell_s}$, where the factor $2$ is absorbed
by lowering the exponent by one (valid since $\ell_s\ge 1$).

\emph{Sparse case $\mu_s<\ell_s$.} Here the relevant overload scale is the additive deviation $\Psifn(\mu_s,\ell_s)$, so we
use a hard threshold. Let $Y_i=\1[r_i>K\Psifn(\mu_s,\ell_s)]$, again nondecreasing in $r_i$. Its expectation is precisely
the probability that the binomial load exceeds the threshold $t=K\Psifn$, which we bound by Bennett's inequality, showing
its exponent is at least $B\ell_s$. There are two sub-regimes governed by $\ell_s/\mu_s$. When $\ell_s/\mu_s=\Oh(1)$ the
deviation $t$ is comparable to $\sqrt{\mu_s\ell_s}$, the Gaussian part of Bennett's exponent applies, and it contributes
$\gtrsim t^2/\mu_s\gtrsim K^2\ell_s$. When $\ell_s/\mu_s\to\infty$ the threshold is dominated by the Poissonian summand
$K\ell_s/\log(e+\ell_s/\mu_s)$ of $\Psifn$, and the Bennett exponent $t\log(t/\mu_s)\gtrsim K\ell_s$: the amplification
$\log(t/\mu_s)\asymp\log(\ell_s/\mu_s)$ is exactly cancelled by the denominator $\log(e+\ell_s/\mu_s)$ built into $\Psifn$,
while the additive ``$1$'' inside $\Psifn$ keeps the bound meaningful as $\mu_s\to0$. In both sub-regimes the exponent is at
least $B\ell_s$ once $K$ is large in terms of $B$, so $\E Y_i=\Pr(r_i>K\Psifn)\le e^{-B\ell_s}$.

\paragraph{From expectation to high probability, and the union.}
We now turn the per-row expectation bound into control of the number of overloaded rows, and then pay for all subsets. In
both cases $Y_i$ is a nondecreasing function of $r_i$, hence so is $e^{tY_i}$ for $t\ge 0$, and negative association lets us
bound the moment generating function of $\sum_i Y_i$ by the product of the individual ones, exactly as under independence;
this is the upper-tail factorization. The count $\sum_i Y_i$ has mean $M=\sum_i\E Y_i\le n e^{-B\ell_s}$, and we ask that it
not exceed $a=s/64$. The standard Chernoff bound for a sum of $\{0,1\}$ (or $[0,1]$) variables with this factorization
gives
\[
  \Pr\Big(\sum_i Y_i>a\Big)\ \le\ \Big(\tfrac{eM}{a}\Big)^{a}\ \le\ \big(64e\,(n/s)\,e^{-B\ell_s}\big)^{s/64}\ \le\ e^{-cB\,s\ell_s},
\]
where the middle step substitutes $M\le n e^{-B\ell_s}$ and $a=s/64$, and the last step uses $\log(n/s)\le\ell_s+\Oh(1)$
(which holds because $T=\Th(n)$, so $\log(n/s)$ and $\ell_s=\log(eT/s)$ differ by $\Oh(1)$); the factor $64e\,(n/s)$ is thus
$e^{\Oh(\ell_s)}$ and is swallowed by the dominant $e^{-Bs\ell_s/64}$, leaving an exponent $cB\,s\ell_s$ for an absolute
constant $c$. This says that for a \emph{fixed} $S$ the probability that more than $s/64$ rows are overloaded is at most
$e^{-cB\,s\ell_s}$. It remains to union over $S$. The number of $s$-subsets is $\binom{T}{s}\le(eT/s)^s=e^{s\ell_s}$, so
summing the failure probability over all subsets of each size $s=1,\dots,T$,
\[
  \Pr(\text{some }S\text{ fails})\ \le\ \sum_{s=1}^{T} e^{s\ell_s}\,e^{-cB\,s\ell_s}\ =\ \sum_{s=1}^{T} e^{-(cB-1)s\ell_s}.
\]
Here the enumeration cost $e^{s\ell_s}$ and the per-subset bound $e^{-cB\,s\ell_s}$ carry the \emph{same} weight $s\ell_s$,
which is what makes a single constant $B$ suffice across all scales. The function $s\ell_s=s\log(eT/s)$ is increasing on
$[1,T]$, so $s\ell_s\ge\ell_1=\log(eT)$ for every $s$; hence each term is at most $(eT)^{-(cB-1)}$, and the sum of at most
$T$ such terms is at most $T(eT)^{-(cB-1)}$. This is $o(1)$ once $cB>2$, which fixes the choice of $B$, and the constant $K$
is the $K(B)$ produced by the tail estimates above. \qed

\subsection{The splitting lemma (\Cref{lem:split})}\label{app:split-full}

The goal of the splitting lemma is to halve a column set $S$ into a part $Y$ of size $a$ so that every row load is split
almost perfectly in proportion, with only an additive $\Psifn$ error. We obtain $Y$ by partial coloring: we run the
Lovett--Meka walk, which finds a sign vector that lands on the boundary of a body cut out by a handful of linear
constraints, and we tune those constraints, one per row, so that the rows whose loads we must control are kept on
near-equality faces while the rest are free to move. The size constraint is enforced by one extra equality. The walk is
then summed over $\Oh(\log m)$ halving stages, and the per-stage errors form a geometric series that collapses to a single
$\Psifn$ term.

Run \Cref{lem:LM} from the starting point $x^{(0)}_t=2a/s-1$ ($t\in S$), so that the target average sign encodes the
intended split size; a sign vector $\sigma\in\{-1,1\}^S$ with $\sum_t\sigma_t=2a-s$ corresponds to the set
$Y=\{t:\sigma_t=1\}$ of size $|Y|=a$. To pin the cardinality we always include the all-ones constraint vector with
threshold $0$, which forces an exact equality $\sum_t\sigma_t^{(\text{step})}=0$ on the increment at every iteration and so
preserves $\sum_t$ throughout the walk; this is what guarantees $|Y|=a$ on exit, not merely in expectation. Let $A$
denote the current alive set of columns, $m=|A|$, and write $\mu_A=dm/n$, $\ell_A=\log(eT/m)$ for the running mean and
logarithmic budget. The Lovett--Meka guarantee is that the walk reaches a face on which at least a constant fraction of the
coordinates are frozen to $\pm1$, provided the constraints leave enough slack; the constraints we impose are exactly the
overload tests of the structural lemma, with locally-tuned thresholds, and we verify the slack in each phase.

\emph{Dense phase ($\mu_A\ge\ell_A$).} In this phase we constrain every nonempty row by a soft Gaussian-width budget scaled
to its own load. For each row with $r_i(A)>0$ set the constraint coefficient $c_i=K\sqrt{\mu_A\ell_A}/\sqrt{r_i(A)}$, so that
heavily loaded rows are held more tightly. The Lovett--Meka feasibility condition is that $\sum_i e^{-c_i^2/16}$ be a small
fraction of $m$, and this is precisely the dense structural estimate: by \eqref{eq:struct-dense}, $\sum_i e^{-c_i^2/16}\le
m/64$. This leaves room for the one extra all-ones equality constraint, since $1+m/64\le m/16$ as soon as $m\ge 22$, so the
total constrained budget stays below the threshold the walk requires. Each iteration moves $\sigma$ along a unit Gaussian
step inside the body, so it changes the load of row $i$ by at most $c_i\|v_i\|_2=c_i\sqrt{r_i(A)}=K\sqrt{\mu_A\ell_A}$,
where the $\sqrt{r_i(A)}$ in $\|v_i\|_2$ cancels the $\sqrt{r_i(A)}$ in the denominator of $c_i$. The per-row, per-stage
change is therefore the uniform bound $K\sqrt{\mu_A\ell_A}$, independent of $i$.

\emph{Sparse phase ($\mu_A<\ell_A$).} Here loads are small and we use a hard threshold. Set $U_A=K\Psifn(\mu_A,\ell_A)$. By
the sparse structural estimate \eqref{eq:struct-sparse}, at most $m/64$ rows have $r_i(A)>U_A$. We constrain exactly those
overloaded rows, with threshold $0$ so the walk does not increase their loads, and we \emph{release} all the remaining rows,
imposing no constraint on them. A released row had at most $U_A$ alive incidences to begin with, so no matter how its
incidences are colored across this and all later stages, its load can change by at most $U_A$ in each of the two halves,
hence by at most $2U_A$ in total. Releasing such rows costs nothing because their entire future deviation is already bounded
by $2U_A=\Oh(\Psifn(\mu_A,\ell_A))$.

\emph{Stage-summing.} It remains to add the per-stage changes along the whole recursion and check that the total is still
$\Oh(\Psifn(\mu_0,\ell_0))$, with no logarithmic blow-up from the number of stages. The alive set halves at each stage, so
after $u$ natural-log steps of halving the mean has shrunk geometrically; writing the shrink factor as $e^{-u}$ we have
$\mu(u)=\mu_0 e^{-u}$ and $\ell(u)=\ell_0+u$, since halving the column count adds a bounded amount to the logarithm at each
stage. For the dense stages, the per-stage change is $K\sqrt{\mu_A\ell_A}=K\sqrt{\mu(u)\ell(u)}=K e^{-u/2}\sqrt{\ell_0+u}$.
The function $e^{-u/2}\sqrt{\ell_0+u}$ is nonincreasing in $u$ for $\ell_0\ge 1$ (the exponential decay dominates the slow
growth of the square root), so summing it over the discrete halving stages is bounded by its value at the start times a
geometric factor, giving $\sum_{\text{dense}}\sqrt{\mu_A\ell_A}\le C\sqrt{\mu_0\ell_0}$. The geometric decay is what
prevents a factor of the number of stages from appearing. For the sparse thresholds, what we need is
$\sup_{\text{sparse}}\Psifn(\mu_A,\ell_A)\le C\Psifn(\mu_0,\ell_0)$, since a released row's contribution is $\Oh$ of a single
such threshold. If $\mu_0\le\ell_0$ then the whole recursion is sparse and $\Psifn(\mu_A,\ell_A)$ is monotone in the stage,
so the supremum is attained at the start and equals $\Psifn(\mu_0,\ell_0)$. If $\mu_0>\ell_0$, the recursion begins dense
and only enters the sparse phase once $\mu(u)$ has dropped below $\ell(u)$, which happens at some stage
$u_*\le\log(\mu_0/\ell_0)$; at that point $\ell_*=\ell_0+u_*\le\ell_0+\log(\mu_0/\ell_0)\le C\sqrt{\mu_0\ell_0}$, so the
sparse phase effectively restarts from a budget already controlled by $\sqrt{\mu_0\ell_0}\le C\Psifn(\mu_0,\ell_0)$, and we
rebase the sparse bound there. The recursion stops once $\Oh(1)$ variables remain alive, at which point the residual
coloring is found by brute force at cost $\Oh(1)$ per row; this constant is absorbed by the additive ``$1$'' inside
$\Psifn$. Finally, the cardinality is exact up to an $\Oh(1)$ cleanup swap to repair any unit imbalance, and such a swap
changes each $r_i(Y)$ by at most $1\le\Psifn$. Dividing the resulting signed discrepancy by $2$ converts the symmetric
$\pm1$ split into the two-sided load bound claimed. \qed

\subsection{The recursive color-tree (\Cref{lem:offline-upper})}\label{app:tree-full}

We build the $k$-coloring by repeatedly halving the column set with \Cref{lem:split}, forming a balanced binary tree whose
leaves are the $k$ color classes. The key point is bookkeeping: each split introduces a fresh additive error into the row
loads, but that error is then carried down to the leaf scaled by the fraction of columns that survive, and these fractions
multiply to a small number. We track one row $i$ and one root-to-leaf path, write the leaf load as the ideal proportional
load plus the accumulated error, and show the accumulated error sums to a single $\Psifn$ with no extra logarithmic factor
from the depth of the tree.

For a node $v$ with $q_v$ leaves below it and column set of size $M_v$, applying \Cref{lem:split} produces children $w$
whose loads satisfy $r_i(w)=(M_w/M_v)r_i(v)+e_i(w)$, where the first term is the perfectly proportional split and the
additive error obeys $|e_i(w)|\le C\Psifn(dM_v/n,\log(eT/M_v))$ by the splitting lemma. We index nodes by the number $q$ of
leaves they subtend. For a $q$-leaf node the relevant parameters are $dM_v/n=\Th(\lam q)$, since such a node holds a
$\Th(q/k)$ fraction of all columns and $\lam=\Th(dk/n)/k$, and $\log(eT/M_v)=\Th(1+\log(k/q))=:\ell_q$. So the error
injected at a $q$-node is $\Oh(\Psifn(\Th(\lam q),\ell_q))$.

Now follow this error to the leaf. Passing the recursion $r_i(w)=(M_w/M_v)r_i(v)+e_i(w)$ down the path expresses the leaf
load $x_{i,c}$ as the ideal load $\tfrac{m_c}{T}\deg(i)$ plus a sum of the injected errors, each multiplied by the product
of the child-fractions $M_{(\cdot)}/M_{(\cdot)}$ on the remaining path from that node to the leaf. Along a balanced path
those fractions are each $\Th(1/2)$ and they telescope: the product of the fractions from a $q$-node down to the leaf is the
ratio of the leaf's column count to the $q$-node's column count, which is $\Th(1/q)$. Hence the error born at a $q$-node
reaches the leaf attenuated by $\Th(1/q)$. Summing the attenuated errors along the root-to-leaf path,
\[
  \Big|x_{i,c}-\tfrac{m_c}{T}\deg(i)\Big|\ \le\ C\sum_{q\text{ on path}}\frac{\Psifn(\Th(\lam q),\ell_q)}{q}
  \ =\ C\Big(\underbrace{\sum_q\tfrac1q}_{\Oh(1)}+\sqrt\lam\sum_q\sqrt{\tfrac{\ell_q}{q}}+\sum_q\frac{\ell_q}{q\,\log(e+\ell_q/(\lam q))}\Big),
\]
where the three groups come from expanding $\Psifn(\Th(\lam q),\ell_q)=\Th(1)+\Th(\sqrt{\lam q\,\ell_q})+\Th(\lam
q\,\ell_q/(\lam q\log(e+\ell_q/(\lam q))))$ and dividing each piece by $q$: the constant ``$1$'' inside $\Psifn$ gives
$\sum_q 1/q$, the Gaussian part $\sqrt{\lam q\,\ell_q}/q=\sqrt\lam\sqrt{\ell_q/q}$ gives the middle sum, and the Bennett part
$(\lam q)\ell_q/(\lam q\cdot\log(e+\ell_q/(\lam q)))$ divided by $q$ gives the last sum after cancelling $\lam q$. The
values of $q$ along the path are the geometric sequence $k,k/2,\dots,2$, so each of the three sums is a geometric-type sum
that we now evaluate.

The first sum is harmless: $\sum_q 1/q$ over the geometric sequence $k,k/2,\dots,2$ is a convergent geometric series, hence
$\Oh(1)$.

\emph{Square-root term.} We must bound $\sum_q\sqrt{\ell_q/q}$ with $\ell_q=1+\log(k/q)$. As $q$ ranges over
$k,k/2,\dots,2$, the factor $1/\sqrt q$ decays geometrically while $\sqrt{\ell_q}$ grows only like $\sqrt{\log(k/q)}$, so the
sum is dominated by its small-$q$ (large-$k/q$) end, where $\ell_q=\Th(\log k)=\Th(\Lscale)$; carrying out the geometric sum
gives $\sum_q\sqrt{\ell_q/q}=\Oh(\sqrt{\Lscale})$. Multiplying by the prefactor $\sqrt\lam$, the whole square-root term is
$\Oh(\sqrt{\lam\Lscale})$.

\emph{Bennett term.} We must bound $\sum_q \ell_q/\big(q\,\log(e+\ell_q/(\lam q))\big)$. If $\lam\ge\Lscale$, then for every
$q\ge1$ we have $\ell_q/(\lam q)\le\Lscale/\lam\le 1$, so the denominator $\log(e+\ell_q/(\lam q))\ge1$, and the sum is at
most $\sum_q\ell_q/q=\Oh(\Lscale)=\Oh(\sqrt{\lam\Lscale})$, the last equality holding because $\lam\ge\Lscale$. If
$\lam<\Lscale$, put $A=\Lscale/\lam>1$ and split the path at $q=\sqrt A$. For the small-$q$ part $q\le\sqrt A$ we have
$\ell_q\asymp\Lscale$ and $\ell_q/(\lam q)\ge\Lscale/(\lam\sqrt A)=\sqrt A$, so $\log(e+\ell_q/(\lam q))\ge c\log(e+A)$, and
the geometric sum $\sum_{q\le\sqrt A}\ell_q/q$ is $\Oh(\Lscale)$, giving an overall $\Oh(\Lscale/\log(e+A))$ for this part.
For the large-$q$ part $q>\sqrt A$ we simply drop the denominator (it is $\ge1$) and bound $\sum_{q>\sqrt A}\ell_q/q\le
C\Lscale/\sqrt A=C\Lscale/\sqrt{\Lscale/\lam}=C\sqrt{\lam\Lscale}$. Adding the two parts, the Bennett term is
$\Oh(\sqrt{\lam\Lscale}+\Lscale/\log(e+\Lscale/\lam))$.

Adding the three pieces, $\Oh(1)+\Oh(\sqrt{\lam\Lscale})+\Oh(\sqrt{\lam\Lscale}+\Lscale/\log(e+\Lscale/\lam))$, recombines
exactly into $\Oh(\Psifn(\lam,\Lscale))$ with $\Lscale=\log(ek)$; no extra logarithmic factor survives because the
attenuation $1/q$ turned each path sum into a geometric series dominated by one endpoint. Finally we pass from per-color
loads to the range. The ideal loads of two colors differ by
$|(m_c/T)\deg(i)-(m_{c'}/T)\deg(i)|\le\deg(i)/T\le 1$, because the class sizes $m_c$ differ by at most one and
$\deg(i)\le T$; therefore $\rangeop_i\le 2\max_c|x_{i,c}-(m_c/T)\deg(i)|+1=\Oh(\Psifn(\lam,\log k))$ at every coordinate
$i$. The whole construction is deterministic once the incidence matrix $A$ is fixed, because the structural event of
\Cref{lem:structural} simultaneously covers every alive set the recursion can select, so no further randomness is spent
beyond drawing $A$. \qed
